\section*{Stone--Weierstrass}

\begin{theorem}[Stone--Weierstrass]
Sia $X$ uno spazio compatto.
Sia $\mathfrak A \subset C(X)$ un'algebra 
cioè un sottospazio vettoriale dello spazio delle funzioni continue 
su $X$ a valori in $\RR$
tale che se $f,g\in \mathfrak A$ allora anche $f\cdot g\in \mathfrak A$.

\begin{itemize}
\item $1 \in \mathfrak A$
\item $\mathfrak A$ separa i punti di $X$, cioè
\[
\forall x \neq y \in X \quad \exists f \in \mathfrak A \text{ tale che } f(x)\neq f(y).
\]
\end{itemize}

Allora $\mathfrak A$ è densa in $C(X)$ rispetto alla norma uniforme:
\[
\overline{\mathfrak A}^{\|\cdot\|_\infty} = C(X)
\]
cioè 
\[
\forall \eps>0 \forall f\in C(X)\exists g\in \mathfrak A\colon \Abs{f-g}_\infty < \eps.
\]
\end{theorem}

Per la dimostrazione ci servono alcuni lemmi intermedi.

% Se $x,y\in X$ con $x\neq y$, per ipotesi esiste
% $f\in\mathfrak A$ tale che
% 
% \[
% f(x)\neq f(y).
% \]
% 
% Riscalando la funzione possiamo assumere
% 
% \[
% f(x)=0, \qquad f(y)=1 .
% \]
% 
% Per continuità esistono intorni $U_x,U_y$ tali che
% 
% \[
% f(U_x) \subset (-\tfrac13,\tfrac13),
% \qquad
% f(U_y) \subset (\tfrac23,\tfrac43).
% \]
% 
% Quindi in $U_x$ la funzione è vicina a $0$ mentre in $U_y$
% è vicina a $1$.

\begin{lemma}
Siano $X$ e $\mathfrak A$ come nelle ipotesi del teorema.
Dato $x_0\in X$, dato $U$ intorno aperto di $x_0$
esiste $V$ intorno di $x_0$, $V\subset U$
tale che per ogni $\eps>0$ 
esiste $f\in \mathfrak A$ tale che 
\begin{gather*}
    0 \le f(x) \le 1\\
    f(x) < \eps\qquad \forall x\in V\\
    f(x) > 1-\eps \qquad \forall x \in X\setminus U
\end{gather*}
\end{lemma}
\begin{proof}
Per ogni $x\in X\setminus U$ 
consideriamo $g_x \in mathfrak A$ 
tale che $g_x(x)\neq g_x(x_0)$.
Definiamo $h_x = g_x - g_x(x_0)\cdot 1 \in \mathfrak A$.
Così $h_x(x) \neq  h_x(x_0) = 0$.
Definiamo
\[
 p_x = \frac{h_x^2}{\Abs{h_x}_\infty} \in \mathfrak A
\]
così $p_x(x_0)=0$, $p_x(x)>0$,
$0\le p_x \le 1$.

L'insieme 
\[
U_x = \ENCLOSE{g\in X\colon p_x(y)>0}
\]
è un intorno aperto di $x$
dunque $X\setminus U$ è un compatto 
ricoperto dagli $U_x$ al variare di $x\in X\setminus U$.

Posso estrarre un sottoricoprimento finito:
dunque esistono $x_1,\dots,x_n\in X\setminus U$ 
tali che $X\setminus U \subset U_{x_1} \cup \dots \cup U_{x_m}$.

Poniamo
\[
  p = \frac{p_{x_1}+ \dots + p_{x_m}}{m} \in \mathfrak A
\]
così da avere $0\le p \le 1$, $p(x_0)=0$,
$p(x)>0$ per ogni $x\in X\setminus U$.

Visto che $X\setminus U$ è compatto 
esiste $\delta = \min p(X\setminus U)>0$ 
e dunque $p(x)\ge \delta > 0$ e $0<\delta\le 1$.

Sia 
\[
  V= \ENCLOSE{x\in X \colon p(x) < \frac \delta 2}.
\]
Visto che $p(x_0)=0$ si ha che $V$ è un intorno di $x_0$,
e visto che 
$p(x) \ge \delta$ su $X\setminus U$
risulta $V\subset U$.

Sia $k\in \NN$ tale che 
\[
 \frac 1 \delta < k \le \frac 1 \delta + 1 
 = \frac{\delta+1}{\delta}
 < \frac 2 \delta
\]
così $1 < k \delta < 2$.

Poniamo
\[
 q_n(x) = \Enclose{1-p^n(x)}^{k^n} \in \mathfrak A
\]
così $0\le q_n \le 1$, $q_n(x_0) = 1$.
Ricordando la disuguaglianza di Bernoulli
\[
  (1+t)^n \ge 1 + nt, \qquad \forall t\ge -1
\]
per ogni $x\in V$ si ha $k p(x)\le k \frac \delta 2 < 1$ 
dunque
\[
q_n(x) \ge 1- k^n p^n(x)
\ge 1 - \enclose{k \frac \delta 2}^n \to 1
\]
ovvero $q_n \to 1$ uniformemente su $V$.
Mentre se $x\in X\setminus U$ 
si ha $kp \ge k\delta > 1$
\begin{align*}
 q_n(x) 
 &= \Enclose{1-p^n(x)}^{k^n}
 \le \frac{1+k^np^n(x)}{k^np^n(x)}
 \Enclose{1-p^n(x)}^{k^n}\\
 &\le \frac{\Enclose{1+p^n(x)}^{k^n}}{k^n p^n(x)}
 \Enclose{1-p^n(x)}^{k^n}\\
 &= \frac{\Enclose{1+p^{2n}(x)}^{k^n}}{k^n p^n(x)}
 \le \frac{1}{k^n p^n(x)}
 \le \frac{1}{(k\delta)^n} \to 0.
\end{align*}
Dunque $q_n\to 0$ uniformemente su $X\setminus U$.

Prendendo $n$ sufficientemente grande 
si ha $q_n<\eps$ su $X\setminus U$ 
e $q_n > 1-\eps$ su $V$.

La funzione $f=1-q_n$ soddisfa le proprietà richieste.
\end{proof}

\begin{lemma}
Siano $X$ e $\mathfrak A$ come nel teorema.
Siano $A,B\subset X$ compatti disgiunti.
Per ogni $\eps>0$ esiste $f\in \mathfrak A$ tale che 
\begin{gather*}
    0 \le f \le 1, \\
    f(x) < \eps \qquad \forall x\in A,\\
    f(x) > 1-\eps \qquad \forall x \in B.
\end{gather*}
\end{lemma}
%
\begin{proof}
Sia $U\supset A$ aperto disgiunto da $B$.
Per ogni $x\in A$ esiste $V_x\subset U$ 
un intorno aperto di $x$ come nel lemma precedente.
Gli aperti $V_x$ ricoprono il compatto $A$ 
dunque esistono $x_1, \dots, x_m\in A$ tali che 
$V_{x_1} \cup \dots \cup V_{x_m} \supset A$
ed esistono $g_1, \dots, g_m \in \mathfrak A$ 
tali che $0\le g_k \le 1$, $g_i < \frac \eps m$ su $V_{x_i}$
e $g_i > 1 - \frac \eps m$ su $X\setminus U$.

Posto $g= g_1 \cdot g_2 \cdots g_m$ 
si ha $0\le g \le 1$ e per ogni $x\in A$ 
esiste $x_i$ tale che $x\in V_{x_i}$.
Visto che $g_i(x) < \frac \eps m$ si ha 
$g(x) < \frac \eps m$.
Inoltre per ogni $x\in B$, $x \not \in \bigcup V_{x_i}$
e per ogni $i$ si ha $g_i(x) < 1 - \frac \eps m$.
Dunque 
\[
g > \enclose{1-\frac\eps m}^m \ge 1 - m \frac \eps m = 1-\eps.
\]
\end{proof}

\begin{proof}[Dimostrazione del teorema.]

Sia $f\in C(X)$ e $\eps>0$.
Eventualmente considerando $f+ \Abs{f}_\infty$ posso 
supporre $f\ge 0$.
Supponiamo $\eps < \frac 1 3$ 
e consideriamo $n$ tale che $(n-1)\eps \ge \Abs{f}_\infty$.

Per $j=0, \dots, n$ definiamo 
\begin{align*}
    A_j &= \ENCLOSE{x\in X\colon f(x) \le \enclose{j-\frac 1 3}\eps}\\
    B_j &= \ENCLOSE{X\in X\colon f(x) \ge \enclose{j+\frac 1 3}\eps}
\end{align*}
così $A_j\cap B_j\neq \emptyset$, $A_j$ e $B_j$ sono chiusi,
\begin{gather*}
    \emptyset \subset A_0 \subset A_1 \subset \dots \subset A_n = X, \\
    B_0 \supset B_1 \supset \dots \supset B_n = \emptyset.
\end{gather*}


Per il lemma precedente esiste
$g_j\in\mathfrak A$ tale che
$g_j \le \frac \epsilon n$
su $A_j$ 
mentre $g_j \ge 1 - \frac \eps n$ su $B_j$.

Si definisce
\[
g(x)=\eps\sum_{i=0}^n g_i(x)
\]
così per ogni $x\in X$ esiste $j$ 
tale che se $x\in A_j \setminus A_{j-1}$
si ha 
\[
 \enclose{j - \frac 4 3} \eps 
 < f(x) 
 \le \enclose{j - \frac 1 3} \eps.
\]
Per ogni $i\le j-2$ se $x\in B_i$ 
si ha $g_i(x) < 1 - \frac \eps n$.

Posto
\[
g(x) = \eps \sum_{i=0}^{j-1} g_i(x)
\]
si ha  
\[
g(x)\le \eps \sum_{i=j}^{n} g_i(x)
\le j\eps + \eps (n-j+1)\frac \eps n
\le j\eps + \eps^2 
< \enclose{j+\frac 1 3}\eps
\]
e
\[
g(x) 
\ge \eps \sum_{i=0}^{j-2} g_i(x) 
\ge (j-1) \eps \enclose{1-\frac \eps n}
= (j-1)\eps - \frac{j-1}{n}\eps^2 
> \dots > \enclose{j-\frac 4 3} \eps.
\]

In conclusione 
\[
  \abs{f(x) - g(x)} 
  \le \enclose{j+\frac 1 3} \eps - \enclose{j-\frac 4 3} \eps 
  < 2 \eps.
\]
\end{proof}